%% ============================================================
%% Refining-grid inference at fixed leaf budget (the "inferential half")
%% Standalone pre-integration draft
%%
%% SUPERSEDED (2026-08-31, later same session): the user decided this
%% material should NOT be a pre-integration fragment for theory.tex, but
%% its own fully standalone document with independently elaborated
%% assumptions -- not an "add-on to the current theory doc." The result
%% is `refining-grid-standalone.tex` (same directory), which restates
%% every assumption and background result in full (not by \ref into
%% theory.tex), fixes several defects found in this fragment across a
%% 4-pass review (proof-auditor, theory-devils-advocate, domain-reviewer,
%% Oracle -- see session_notes/doubletree-2026-08-31.md), and is the
%% current, reviewed artifact. This file is kept only for its drafting
%% history and the audit trail of what changed and why; do not integrate
%% it directly, and do not treat it as reflecting the corrected content.
%% Intended location: immediately after _refining-grid-draft.tex's
%% Assumption \ref{ass:mesh} and Lemma \ref{lem:mesh-approx} (the KEPT
%% population-level half), still within the "Robustness to grid
%% alignment" material planned for Sec.~\ref{sec:setup}/Sec.~3.
%%
%% PROVENANCE (2026-08-31 pick-up session, per
%%   quality_reports/plans/2026-08-31_refining-grid-inferential-half-respec.md):
%% This spec explicitly deferred drafting until two gates were re-run
%% against the current (post-"Slice 0") theory.tex:
%%   (1) a fresh proof-auditor pass on Defect 1 (grid-finiteness conflict),
%%       since the original 2026-08-28 audit read the OLD ass:finite;
%%   (2) an Oracle consultation settling whether the "efficiency
%%       corollary" removed on 2026-08-28 is actually a corollary of
%%       thm:anchor/cor:width or an independent theorem.
%% Both gates were run this session. Verdicts, in brief (full detail in
%% the two background-task transcripts referenced by the session note):
%%   - Defect 1 resolves as "genuine dependency found," not "silently
%%     resolved by reading at fixed row n." Four dependencies: D1 (the
%%     population margin Delta_j -> 0 under refinement; UNREPAIRABLE by
%%     any uniformity assumption in the target regime h_n=o(n^{-1/2}),
%%     since rescuing the margin route would need
%%     Lbar*log(d/h_n)=o(n*h_n^2), whose RHS->0 while LHS diverges), D2
%%     (union-bound cardinality |cT_{Lbar,n}| diverges -- repairable via
%%     U1 below), D3 (atom-mass floor min_a P(a) -> 0 -- repairable via
%%     U2 below, relocating the floor to comparator leaves), D4 (CRITICAL,
%%     distinct from D1-D3: ass:sparsity is FALSE at every finite row n,
%%     not just asymptotically, so the entire sparsity-dependent machinery
%%     -- lem:invariance, lem:margin, lem:selection, lem:U, lem:cross,
%%     lem:uniform, thm:main, cor:variance, prop:rate-verify -- is simply
%%     inapplicable at any row; any new result must enter through
%%     thm:crossfit/thm:anchor/cor:width instead). A fifth finding, D5 (an
%%     incorrect Boundaries paragraph contradicting this file's own
%%     TODO(refining-grid) comments, plus two vestigial ass:finite
%%     citations in thm:crossfit/lem:invariance), was fixed directly in
%%     theory.tex this same session, independent of the material below.
%%   - The framing question resolved as a SPLIT verdict, not one corollary:
%%     the cross-fitted CLT is a true corollary of thm:crossfit (NOT of
%%     thm:anchor/cor:width); the anchor-interval width is a corollary of
%%     cor:width but delivers RATE, not EFFICIENCY; the full-sample
%%     estimator's CLT is an INDEPENDENT THEOREM requiring off-sparsity,
%%     vanishing-misspecification replacements for lem:U/lem:cross/
%%     lem:uniform/lem:selection's selection-event device, and a new
%%     cor:variance limit argument.
%% User's explicit choice (2026-08-31, after reviewing both verdicts):
%% full scope -- draft BOTH the true-corollary result and the
%% independent-theorem result, not just the safer corollary-only subset.
%% Drafted by a proof-writer subagent seeded with both verdicts verbatim;
%% NOT yet reviewed by proof-auditor/theory-devils-advocate/domain-reviewer/
%% Oracle (that Phase-3-style multi-agent review, per the spec's Step 4,
%% is deliberately deferred until after a first user check-in on this
%% draft -- see "Compatibility and open items" at the tail of this file
%% for the drafting agent's own confidence assessment of which parts are
%% least likely to survive that review unchanged).
%%
%% Reuses (all pre-existing, either in theory.tex or in the KEPT
%% _refining-grid-draft.tex): ass:mesh, lem:mesh-approx, eq:mesh-slab,
%% eq:mesh-en, eq:mesh-mun (this file); ass:iid, ass:causal, ass:moments,
%% ass:construct, ass:global, ass:rate, thm:crossfit, thm:anchor,
%% cor:width, lem:skeleton, prop:bilinear, lem:biasbound, cor:convert,
%% lem:normequiv, lem:clip, lem:tangent, thm:regular, eq:interior,
%% eq:vhat, eq:feasible, app:honest (theory.tex).
%% ============================================================

\subsection{Inference under grid refinement}
\label{sec:grid-inference}

Continue to use the continuum covariate $Z$, discretisation
$X_n=Q_n(Z)$, grids $\cA_n$, candidate classes
$\cT_{\bar L,n}$, mesh $h_n$, and discretised population nuisances
$\ezn,\mzn$ from Assumption~\ref{ass:mesh}.  The observed data for causal
identification remain $(Z,A,Y)$; $X_n$ is the analyst's row-$n$ feature map
used to construct the tree fits.  Thus the true nuisances in the influence
function remain
\[
  \ez^\circ(Z)=P(A=1\mid Z),
  \qquad
  \mz^\circ(Z)=E(Y\mid A=0,Z),
\]
whereas $\ezn$ and $\mzn$ are their population projections onto
$\sigma(X_n)$.  This distinction is essential: ignorability given $Z$ does
not, in general, imply ignorability given the coarsening $X_n$.

For $j=e$, let
\[
  \|f\|_e^2=E\{f(Z)^2\},
\]
and, for $j=\mu$, let
\[
  \|f\|_\mu^2
  =E[\{1-\ez^\circ(Z)\}f(Z)^2].
\]
For a row-$n$ fitted nuisance define
\[
  D_{e,n}=\|\widehat e_n-\ez^\circ\|_e,
  \qquad
  D_{\mu,n}=\|\widehat\mu_n-\mz^\circ\|_\mu,
  \qquad
  D_{w,n}=\|\widehat w_n-w_0^\circ\|_\mu,
\]
where $w_0^\circ=\ez^\circ/(1-\ez^\circ)$ and
$\widehat w_n=\widehat e_n/(1-\widehat e_n)$ after the one-sided clipping of
Assumption~\ref{ass:construct}(c).

Put
\[
  s_n:=1+\log\card{\cT_{\bar L,n}},
  \qquad
  r_n:=\sqrt{\frac{s_n}{n}}.
\]

\begin{assumption}[Uniform refining-grid inferential conditions]
\label{ass:grid-inference}
In addition to Assumption~\ref{ass:mesh}, the following conditions hold.

\begin{enumerate}[label=(U\arabic*),leftmargin=*]

\item\label{u:entropy}
There is a constant $C_T<\infty$, independent of $n$, such that
\[
  \log\card{\cT_{\bar L,n}}
  \le C_T\bar L\log(dJ_{n,\max}),
  \qquad
  s_n=o(n).
\]
For the root-$n$ conclusions below we impose the stronger condition
\[
  s_n=o(\sqrt n).
  \tag{U1$^+$}\label{eq:u1plus}
\]
When $J_{n,\max}\asymp h_n^{-1}$ and $h_n$ is polynomial in $n$,
$s_n=O\{\bar L\log(d/h_n)\}=O(\log n)$, so both conditions hold.

\item\label{u:comparator}
For each $j\in\{e,\mu\}$, let $\tau^{\mathrm{comp}}_{j,n}$ be the grid-rounded
tree constructed in the proof of Lemma~\ref{lem:mesh-approx}.  After deleting
$P_Z$-null continuum leaves, there is a constant $q_{\min}>0$, independent of
$n$, such that, for all sufficiently large $n$,
\[
  \min_{\ell\in\tau^{\mathrm{comp}}_{j,n}}P\{X_n\in\ell\}
  \ge q_{\min}.
\]
Moreover $m_n/n\to0$, so these comparator trees belong to the feasible sets
in~\eqref{eq:feasible} with probability tending to one.

\item\label{u:bernstein}
The squared-loss fits are uniformly bounded: propensity fits lie in $[0,1]$
before and after clipping, and outcome fits lie in $[-B_Y,B_Y]$.  The penalty
satisfies
\[
  \bar L\lambda_n=O(r_n^2).
\]
For each nuisance and each training sample of order $n$, the finite union of
at-most-$\bar L$-dimensional leaf-constant models satisfies the localized
quadratic-loss bound
\begin{equation}\label{eq:grid-bernstein}
 \sup_{f\in\mathcal F_{j,n}}
 \frac{\left|(P_n-P)\{\ell_j(f)-\ell_j(\gamma_{0,j}^\circ)\}\right|}
      {\|f-\gamma_{0,j}^\circ\|_j r_n+r_n^2}
 =O_p(1),
\end{equation}
where
\[
 \ell_e(f)=(A-f(X_n))^2,\qquad
 \ell_\mu(f)=(1-A)\{Y-f(X_n)\}^2,
\]
and $\mathcal F_{j,n}$ is the union, over
$\tau\in\cT_{\bar L,n}$, of the bounded functions constant on the leaves of
$\tau$.

\item\label{u:approx}
The grid-alignment errors satisfy
\[
  \delta^*_{j,n}=O(\sqrt{h_n}),
  \qquad j\in\{e,\mu\},
\]
as established by Lemma~\ref{lem:mesh-approx}.

\end{enumerate}
\end{assumption}

\paragraph{Why (U2) is derived rather than postulated.}
Let $\ell^\circ$ be a positive-mass leaf of one of the fixed continuum trees
$\tau_j^\circ$, and let $\ell_n$ be its rounded counterpart.  The proof of
Lemma~\ref{lem:mesh-approx} constructs a disagreement set $B_n$ satisfying
\[
  P_Z(B_n)\le(\bar L-1)C_Xh_n.
\]
Since $\ell_n\triangle\ell^\circ\subseteq B_n$,
\[
  |P_Z(\ell_n)-P_Z(\ell^\circ)|
  \le P_Z(\ell_n\triangle\ell^\circ)
  \le(\bar L-1)C_Xh_n.
\]
There are finitely many positive-mass continuum leaves.  Hence
\[
  q^\circ_{\min}
  :=\min_{j\in\{e,\mu\}}
    \min_{\ell^\circ\in\tau_j^\circ:P_Z(\ell^\circ)>0}
       P_Z(\ell^\circ)>0.
\]
For all sufficiently large $n$,
\[
  P_Z(\ell_n)\ge q^\circ_{\min}/2.
\]
Thus (U2) holds with $q_{\min}=q^\circ_{\min}/2$.  By
$1-\ez^\circ\ge c$, every such leaf has control mass at least
$cq_{\min}$.  Finally, $m_n/n\to0$ and the weak law of large numbers imply
that each comparator leaf eventually contains at least $m_n$ observations
(and at least $m_n$ controls for the outcome fit) with probability tending
to one.

\begin{lemma}[Refining-grid oracle inequality]
\label{lem:grid-oracle}
Suppose Assumptions~\ref{ass:iid}, \ref{ass:causal},
\ref{ass:construct}, \ref{ass:global}, \ref{ass:mesh}, and
\ref{ass:grid-inference} hold.  Let $\widehat\gamma_{j,n}$ be the penalized
empirical-risk minimizer over the row-$n$ feasible class.  Then, for each
$j\in\{e,\mu\}$,
\begin{equation}\label{eq:grid-oracle}
  \|\widehat\gamma_{j,n}-\gamma_{0,j}^\circ\|_j
  =O_p\!\left(\sqrt{h_n}+r_n\right).
\end{equation}
More generally, if a population oracle attaining $\delta^*_{j,n}$ itself
has the uniform leaf-mass property in (U2), then
\[
  \|\widehat\gamma_{j,n}-\gamma_{0,j}^\circ\|_j
  =O_p(\delta^*_{j,n}+r_n).
\]
After clipping,
\begin{equation}\label{eq:grid-weight-oracle}
  D_{w,n}=O_p(\sqrt{h_n}+r_n).
\end{equation}
\end{lemma}

\begin{proof}
\emph{Proof strategy.}
Use the grid-rounded tree as an actually feasible comparator.  Apply the
penalized ERM basic inequality, control its empirical-process term by the
localized Bernstein bound~\eqref{eq:grid-bernstein}, and absorb the resulting
linear term into the quadratic population excess risk.  This argument compares
risks and never compares a shrinking population margin.

Fix $j$.  Write $g_j=\gamma_{0,j}^\circ$,
$\widehat f=\widehat\gamma_{j,n}$, and let $f^{\mathrm{comp}}_{j,n}$ be the
population leaf projection on $\tau^{\mathrm{comp}}_{j,n}$.  By (U2), the
comparator is feasible with probability tending to one.  On that event,
optimality in~\eqref{eq:select} gives
\begin{align}
 P_n\ell_j(\widehat f)+\lambda_n\card{\widehat\tau_j}
 &\le
 P_n\ell_j(f^{\mathrm{comp}}_{j,n})
   +\lambda_n\card{\tau^{\mathrm{comp}}_{j,n}},\\
 P_n\{\ell_j(\widehat f)-\ell_j(g_j)\}
 &\le
 P_n\{\ell_j(f^{\mathrm{comp}}_{j,n})-\ell_j(g_j)\}
   +\bar L\lambda_n.
 \label{eq:grid-basic}
\end{align}

For squared loss and a conditional-mean truth, the population excess-risk
identity is exact.  For the propensity,
\[
 P\{\ell_e(f)-\ell_e(\ez^\circ)\}
 =E\{f(X_n)-\ez^\circ(Z)\}^2=\|f-\ez^\circ\|_e^2.
\]
Indeed, expanding the squares gives
\[
 (A-f)^2-(A-\ez^\circ)^2
 =(f-\ez^\circ)^2-2(A-\ez^\circ)(f-\ez^\circ),
\]
and the second term has expectation zero because
$E(A-\ez^\circ\mid Z)=0$ and $f(X_n)$ is $\sigma(Z)$-measurable.  Similarly,
\[
 P\{\ell_\mu(f)-\ell_\mu(\mz^\circ)\}
 =E[\{1-\ez^\circ(Z)\}\{f(X_n)-\mz^\circ(Z)\}^2]
 =\|f-\mz^\circ\|_\mu^2,
\]
because
$E[(1-A)\{Y-\mz^\circ(Z)\}\mid Z]=0$.

Let
\[
  d=\|\widehat f-g_j\|_j,\qquad
  d_c=\|f^{\mathrm{comp}}_{j,n}-g_j\|_j.
\]
Adding and subtracting population risks in~\eqref{eq:grid-basic}, and applying
\eqref{eq:grid-bernstein} to $\widehat f$ and the comparator, yields
\[
 d^2
 \le d_c^2+O_p(dr_n+r_n^2)+O_p(d_cr_n+r_n^2)
       +\bar L\lambda_n.
\]
By (U3), $\bar L\lambda_n=O(r_n^2)$.  Young's inequality,
$ab\le a^2/4+b^2$, gives
\[
 O_p(dr_n)\le \frac14d^2+O_p(r_n^2),
 \qquad
 O_p(d_cr_n)\le \frac14d_c^2+O_p(r_n^2).
\]
Consequently,
\[
 \frac34d^2\le\frac54d_c^2+O_p(r_n^2),
\]
and therefore
\[
 d^2=O_p(d_c^2+r_n^2),
 \qquad
 d=O_p(d_c+r_n).
\]
The comparator construction in Lemma~\ref{lem:mesh-approx} gives
$d_c=O(\sqrt{h_n})$, proving~\eqref{eq:grid-oracle}.  If an actual minimizer
attaining $\delta^*_{j,n}$ satisfies (U2), use it as the comparator to obtain
the sharper displayed statement.

Finally, on $[0,1-c]$ the map $u\mapsto u/(1-u)$ has derivative at most
$c^{-2}$.  The mean-value theorem, Lemma~\ref{lem:clip}, and
Lemma~\ref{lem:normequiv} therefore give
\[
 D_{w,n}
 \le c^{-2}\|\widehat e_n-\ez^\circ\|_\mu
 \le c^{-2}\|\widehat e_n-\ez^\circ\|_e
 =O_p(\sqrt{h_n}+r_n).
\]
\end{proof}

\paragraph{Post-proof audit.}
\emph{Assumptions:} boundedness supplies the Bernstein envelope; (U1) controls
the number of partition models; (U2) makes the comparator feasible; (U3)
controls both local stochastic fluctuations and the penalty; and (U4) supplies
its approximation rate.
\emph{Dependencies:} only the squared-loss excess-risk identity,
\eqref{eq:grid-bernstein}, Young's inequality,
Lemma~\ref{lem:mesh-approx}, Lemma~\ref{lem:clip}, and
Lemma~\ref{lem:normequiv} are used.
\emph{Rates:} the quadratic inequality yields norm error
$O_p(\sqrt{h_n}+r_n)$, not $O_p(\sqrt{h_n}+n^{-1/2})$ unless
$s_n=O(1)$.
\emph{Edge cases:} the result can fail if the comparator loses leaf mass, the
outcome is unbounded without a replacement tail condition, or the penalty
dominates $r_n^2$.
\emph{Tightness:} the $\sqrt{s_n/n}$ model-selection term is generally
unavoidable over a growing finite class.
\emph{Failure localization:} the substantive condition is
\eqref{eq:grid-bernstein}; it is the Bernstein/curvature replacement for the
invalid shrinking-margin argument.
No object is treated as invariant across $n$: the grid, class, comparator,
projections, and fitted functions all carry an explicit row index.

\begin{corollary}[Cross-fitted inference on a refining grid]
\label{cor:grid-crossfit}
Suppose the conditions of Lemma~\ref{lem:grid-oracle} hold for every fold-wise
training sample, with fixed $K$.  Suppose additionally that
\[
  h_n=o(n^{-1/2}),
  \qquad
  s_n=o(\sqrt n).
  \label{eq:grid-rootn}
\]
Let $\thetatil_n$ be the cross-fitted estimator of
Theorem~\ref{thm:crossfit}, constructed from the row-$n$ grids.  Then
\[
  \sqrt n(\thetatil_n-\theta_0)\rightsquigarrow N(0,V).
\]
\end{corollary}

\begin{proof}
\emph{Proof strategy.}
Verify Assumption~\ref{ass:rate} using Lemma~\ref{lem:grid-oracle}, and then
apply Theorem~\ref{thm:crossfit} verbatim.  Cross-fitting removes any need for
uniform empirical-process control at the score-evaluation stage.

For each fixed fold, its training sample has size
$n(1-1/K)+O(1)\asymp n$, so Lemma~\ref{lem:grid-oracle} gives
\[
 \|\widehat e_n-\ez^\circ\|_e
   =O_p(\sqrt{h_n}+r_n),
 \qquad
 \|\widehat\mu_n-\mz^\circ\|_\mu
   =O_p(\sqrt{h_n}+r_n).
\]
Their product is
\begin{align}
 \|\widehat e_n-\ez^\circ\|_e
 \|\widehat\mu_n-\mz^\circ\|_\mu
 &=O_p\{(\sqrt{h_n}+r_n)^2\}\\
 &=O_p\left(h_n+2\sqrt{h_n}\,r_n+r_n^2\right).
\end{align}
Multiplying each deterministic rate by $\sqrt n$ gives
\[
 \sqrt n\,h_n\to0,
 \qquad
 \sqrt n\,r_n^2=\frac{s_n}{\sqrt n}\to0,
\]
by~\eqref{eq:grid-rootn}, and
\[
 \sqrt n\,\sqrt{h_n}\,r_n
 =\sqrt{h_ns_n}
 =\sqrt{(\sqrt n h_n)(s_n/\sqrt n)}
 \longrightarrow0.
\]
Hence
\[
 \|\widehat e_n-\ez^\circ\|_e
 \|\widehat\mu_n-\mz^\circ\|_\mu
 =o_p(n^{-1/2}),
\]
which is precisely Assumption~\ref{ass:rate}, with the covariate in the score
understood to be $Z$ and the fitted nuisances restricted to functions of
$X_n=Q_n(Z)$.

The remaining hypotheses of Theorem~\ref{thm:crossfit} are
Assumptions~\ref{ass:iid}--\ref{ass:moments},
\ref{ass:construct}, and \ref{ass:global}; they are included above.
Its proof is grid-free after conditioning on each training fit: the held-out
empirical-process term is a centred average independent of that fit.
Theorem~\ref{thm:crossfit} therefore yields
\[
 \sqrt n(\thetatil_n-\theta_0)\rightsquigarrow N(0,V).
\]
\end{proof}

\paragraph{Post-proof audit.}
\emph{Assumptions:} the stronger entropy requirement
$s_n=o(\sqrt n)$ is necessary because the oracle norm error contains
$\sqrt{s_n/n}$.
\emph{Dependencies:} Lemma~\ref{lem:grid-oracle} verifies
Assumption~\ref{ass:rate}; Theorem~\ref{thm:crossfit} then applies directly.
\emph{Rates:} all three terms in
$(\sqrt{h_n}+r_n)^2$ are explicitly $o(n^{-1/2})$.
\emph{Edge cases:} U1 alone, $s_n=o(n)$, is insufficient; it permits
$s_n/\sqrt n\nrightarrow0$.
\emph{Tightness:} the mesh condition is sufficient and matches the product
of two $O(\sqrt{h_n})$ approximation errors.
\emph{Failure localization:} failure can occur at the DML product-rate step,
not in Theorem~\ref{thm:crossfit}.
Nothing is assumed literally invariant across rows except the fixed constants
$d,\bar L,c,B_Y,C_T$, and $q_{\min}$.

\begin{corollary}[Anchor-interval width under grid refinement]
\label{cor:grid-width}
Suppose the conditions of Corollary~\ref{cor:grid-crossfit} hold and the
cross-fitted variance estimator satisfies
$\widetilde\sigma_n^2\to_pV$.  Form the anchor interval of
Theorem~\ref{thm:anchor} using $\thetatil_n$.  Then
\[
 \liminf_{n\to\infty}P(\theta_0\in C_n)\ge1-\alpha
\]
and
\begin{equation}\label{eq:grid-width}
 \operatorname{width}(C_n)
 =O_p\!\left(
   \max\left\{
     \sqrt{\frac{\bar L\log(dJ_{n,\max})}{n}},
     D_{w,n}D_{\mu,n}
   \right\}
 \right).
\end{equation}
Under Lemma~\ref{lem:grid-oracle},
\[
 D_{w,n}D_{\mu,n}
 =O_p\left(h_n+\sqrt{h_n}\,r_n+r_n^2\right).
\]
This is a width-rate statement, not an efficiency statement.
\end{corollary}

\begin{proof}
\emph{Proof strategy.}
Use Theorem~\ref{thm:anchor} for validity.  Repeat the decomposition in the
proof of Corollary~\ref{cor:width}, replacing its fixed-class union bound by
the growing-class rate $r_n$.

Theorem~\ref{thm:anchor} applies because
Corollary~\ref{cor:grid-crossfit} supplies a valid asymptotically normal anchor
and $\widetilde\sigma_n^2\to_pV>0$.  Thus validity follows without any
condition on the full-sample tree estimator.

For width, write
\[
 \widehat\delta_n
 =(\widehat\theta_n-\theta_0)-(\thetatil_n-\theta_0).
\]
By Corollary~\ref{cor:grid-crossfit},
$\thetatil_n-\theta_0=O_p(n^{-1/2})$.  Proposition~\ref{prop:bilinear} and
Lemma~\ref{lem:biasbound} give
\[
 |\theta(\widehat\eta_n)-\theta_0|
 \le\pi^{-1}D_{w,n}D_{\mu,n}.
\]
The fixed-class $O_p(n^{-1/2})$ plug-in fluctuation used in
Corollary~\ref{cor:width} becomes $O_p(r_n)$ under (U1) and
\eqref{eq:grid-bernstein}.  Therefore
\[
 |\widehat\theta_n-\theta_0|
 =O_p\{D_{w,n}D_{\mu,n}+r_n\},
\]
and hence
\[
 |\widehat\delta_n|
 =O_p\{D_{w,n}D_{\mu,n}+r_n\}.
\]
Because
\[
 \operatorname{width}(C_n)
 =2\left\{|\widehat\delta_n|
   +z_{1-\alpha/2}\widetilde\sigma_n/\sqrt n\right\},
\]
and $r_n\ge n^{-1/2}$ by definition of $s_n$, this proves
\[
 \operatorname{width}(C_n)
 =O_p\{\max(r_n,D_{w,n}D_{\mu,n})\}.
\]
Finally, Lemma~\ref{lem:grid-oracle} gives
\[
 D_{w,n}D_{\mu,n}
 =O_p\{(\sqrt{h_n}+r_n)^2\}
 =O_p(h_n+\sqrt{h_n}\,r_n+r_n^2).
\]
Substitution and (U1) yield~\eqref{eq:grid-width}.
\end{proof}

\paragraph{Post-proof audit.}
\emph{Assumptions:} anchor validity and width control are logically separate;
only the latter uses the oracle rates.
\emph{Dependencies:} Theorem~\ref{thm:anchor},
Corollary~\ref{cor:grid-crossfit},
Proposition~\ref{prop:bilinear}, and Lemma~\ref{lem:biasbound} are used.
\emph{Rates:} the fixed-class $n^{-1/2}$ term is replaced by
$r_n=\sqrt{s_n/n}$.
\emph{Edge cases:} coverage survives even when the width fails to contract,
provided the anchor remains valid.
\emph{Tightness:} no matching lower bound for the width is asserted.
\emph{Failure localization:} this corollary does not establish
$\sqrt n(\widehat\theta_n-\thetatil_n)=o_p(1)$ and therefore does not establish
efficiency of the displayed-tree estimator.
Every class-dependent rate explicitly retains its $n$ index.

\begin{theorem}[Full-sample inference under vanishing grid misfit]
\label{thm:grid-fullsample}
Suppose the conditions of Corollary~\ref{cor:grid-crossfit} hold.  Suppose,
in addition, that the localized Bernstein control
\eqref{eq:grid-bernstein} applies jointly to the following finite unions of
bounded functions:
\[
 \{h(\Pi^\nu_\tau\mzn-\mz^\circ):\tau\in\cT_{\bar L,n}\},
 \quad
 \{(1-A)(w_\tau-w_0^\circ)^2:\tau\in\cT_{\bar L,n}\},
\]
and their pairwise products, where
$h=w_0^\circ(1-A)-A$.  Then the full-sample estimator
$\widehat\theta_n$ formed from the two selected row-$n$ trees satisfies
\[
 \sqrt n(\widehat\theta_n-\theta_0)
 =\sqrt n\,P_n\psi_0+o_p(1)
 \rightsquigarrow N(0,V).
\]
Moreover, the plug-in estimator~\eqref{eq:vhat} satisfies
\[
 \widehat V_n\to_pV,
\]
and the corresponding Wald interval has pointwise asymptotic coverage
$1-\alpha$.

If, additionally, the interiority condition~\eqref{eq:interior} holds, the
contiguity argument of Theorem~\ref{thm:regular} applies to the displayed
asymptotic-linear expansion; thus $\widehat\theta_n$ is regular and locally
asymptotically efficient at the fixed continuum law $P_0$.
\end{theorem}

\begin{proof}
\emph{Proof strategy.}
Start from the exact skeleton identity of Lemma~\ref{lem:skeleton}.  Replace
the sparsity-dependent proofs of Lemmas~\ref{lem:U} and~\ref{lem:cross} by
localized bounds in the vanishing realised errors.  The genuinely new step is
uniformity over the growing, data-selected partition class.  After proving
the two remainders are $o_p(n^{-1/2})$, apply the ordinary central limit
theorem.  Finally prove variance consistency directly from vanishing
empirical $L_2$ error; the fixed-misspecification counterexample in
Appendix~\ref{app:honest} disappears because both discrepancies now vanish.

Let
\[
 a_{\mu,\tau}
   :=\Pi^\nu_\tau\mzn-\mz^\circ,
 \qquad
 v_{\mu,\tau}
   :=\widehat\mu_\tau-\Pi^\nu_\tau\mzn.
\]
Thus
\[
 \widehat\mu_\tau-\mz^\circ=a_{\mu,\tau}+v_{\mu,\tau}.
\]

\emph{Step 1: replacement for Lemma~\ref{lem:U}.}
For fixed $\tau$, decompose
\[
 U(\tau)
 =P_n[h\,a_{\mu,\tau}]
  +P_n[h\,v_{\mu,\tau}].
\]
Since $E(h\mid Z)=0$ by Lemma~\ref{lem:skeleton},
$P[h\,a_{\mu,\tau}]=0$.  The localized Bernstein bound and a peeling argument
over the deterministic class $\cT_{\bar L,n}$ give, simultaneously for all
$\tau$,
\[
 |P_n[h\,a_{\mu,\tau}]|
 =O_p\{\|a_{\mu,\tau}\|_\mu r_n+r_n^2\}.
 \label{eq:grid-U-approx}
\]
For the leaf-refit component, write
$v_{\mu,\tau}=\sum_{\ell\in\tau}v_{\ell,\tau}\mathbf1_\ell$.  Then
\[
 P_n[h\,v_{\mu,\tau}]
 =\sum_{\ell\in\tau}
   v_{\ell,\tau}P_n[h\mathbf1_\ell].
\]
Uniform finite-dimensional least-squares concentration over the same class
gives
\[
 \sum_{\ell\in\tau}
 |v_{\ell,\tau}P_n[h\mathbf1_\ell]|
 =O_p(r_n^2)
 \label{eq:grid-U-refit}
\]
uniformly over $\tau$.  This is the growing-class analogue of the
$O_p(\bar L/n)$ calculation in Lemma~\ref{lem:U}; it uses neither
leaf-constancy of $\mz^\circ$ nor Assumption~\ref{ass:sparsity}.

Evaluating~\eqref{eq:grid-U-approx}--\eqref{eq:grid-U-refit} at the selected
partition and using Lemma~\ref{lem:grid-oracle},
\[
 |U|
 =O_p\{D_{\mu,n}r_n+r_n^2\}.
\]
Consequently,
\[
 \sqrt n\,|U|
 =O_p\{D_{\mu,n}\sqrt{s_n}+s_n/\sqrt n\}.
\]
Now
\[
 D_{\mu,n}\sqrt{s_n}
 =O_p\{\sqrt{h_ns_n}+s_n/\sqrt n\}=o_p(1)
\]
by~\eqref{eq:grid-rootn}.  Hence $\sqrt n\,U=o_p(1)$.

\emph{Step 2: replacement for Lemma~\ref{lem:cross}.}
Steps 1 and 3 of the proof of Lemma~\ref{lem:cross} require neither
sufficiency nor a fixed grid.  For each fixed pair
$(\tau_e,\tau_\mu)$, centre
$\Delta w_\tau=\widehat w_{\tau_e}-w_0^\circ$ within the control observations
of each $\tau_\mu$ leaf:
\[
 b_i=\Delta w_\tau(Z_i)
 -\overline{\Delta w}_{\ell_\mu(i)}.
\]
The exact refit identity gives
\[
 T(\tau_e,\tau_\mu)
 =n^{-1}\sum_{i:A_i=0}b_i\{Y_i-\widehat\mu_{\tau_\mu}(X_{n,i})\}.
\]
Insert
\[
 Y_i-\widehat\mu_{\tau_\mu}
 =\varepsilon_i-a_{\mu,\tau_\mu}(Z_i)-v_{\mu,\tau_\mu}(X_{n,i}),
 \qquad
 \varepsilon_i=Y_i-\mz^\circ(Z_i).
\]
The term containing $v_{\mu,\tau_\mu}$ vanishes exactly because it is constant
within each $\tau_\mu$ leaf and $\sum_{\ell,A_i=0}b_i=0$.  Therefore
\[
 T=T_\varepsilon-T_a,
\]
where
\[
 T_\varepsilon=n^{-1}\sum_{i:A_i=0}b_i\varepsilon_i,
 \qquad
 T_a=n^{-1}\sum_{i:A_i=0}b_i a_{\mu,\tau_\mu}(Z_i).
\]

Conditionally on $(Z_{1:n},A_{1:n})$, the coefficients $b_i$ are measurable,
the $\varepsilon_i$ are independent and mean zero, and
$E(\varepsilon_i^2\mid Z_i,A_i=0)\le B_Y^2$.  Hence
\[
 \operatorname{Var}(T_\varepsilon\mid Z,A)
 \le B_Y^2n^{-2}\sum_{i:A_i=0}b_i^2.
\]
Centring is an $\ell_2$ projection, so
\[
 \sum_{i:A_i=0}b_i^2
 \le\sum_{i:A_i=0}\Delta w_\tau(Z_i)^2.
\]
The assumed localized empirical-norm control gives, uniformly over
$\tau_e$,
\[
 n^{-1}\sum_{i:A_i=0}\Delta w_\tau(Z_i)^2
 =O_p(D_{w,n}^2+r_n^2).
\]
Conditional Chebyshev therefore yields
\[
 |T_\varepsilon|
 =O_p\{(D_{w,n}+r_n)/\sqrt n\}.
\]

For $T_a$, Cauchy--Schwarz for the empirical control measure gives
\[
 |T_a|
 \le
 \left\{P_n[(1-A)b^2]\right\}^{1/2}
 \left\{P_n[(1-A)a_{\mu,\tau_\mu}^2]\right\}^{1/2}.
\]
The same localized norm comparison, applied jointly over the pair of
partition classes, gives
\[
 |T_a|
 =O_p\{D_{w,n}D_{\mu,n}
       +(D_{w,n}+D_{\mu,n})r_n+r_n^2\}.
\]
Thus
\[
 |T|
 =O_p\!\left[
 \frac{D_{w,n}+r_n}{\sqrt n}
 +D_{w,n}D_{\mu,n}
 +(D_{w,n}+D_{\mu,n})r_n+r_n^2
 \right].
\]
After multiplication by $\sqrt n$, the first term is $o_p(1)$ because
$D_{w,n}+r_n=o_p(1)$.  The product term is $o_p(1)$ by the calculation in
Corollary~\ref{cor:grid-crossfit}.  Moreover,
\[
 \sqrt n(D_{w,n}+D_{\mu,n})r_n
 =O_p\{\sqrt{h_ns_n}+s_n/\sqrt n\}=o_p(1),
\]
and $\sqrt n r_n^2=s_n/\sqrt n\to0$.  Therefore
$\sqrt nT=o_p(1)$.

\emph{Step 3: full-sample asymptotic linearity.}
Lemma~\ref{lem:skeleton} is an exact algebraic identity and does not require
Assumption~\ref{ass:sparsity}.  Hence
\[
 \widehat\pi(\widehat\theta_n-\theta_0)
 =\pi P_n\psi_0+U-T.
\]
Steps 1 and 2 give $\sqrt n(U-T)=o_p(1)$.  Also
$\widehat\pi-\pi=O_p(n^{-1/2})$ and
$\sqrt nP_n\psi_0=O_p(1)$, so, exactly as in the proof of
Lemma~\ref{lem:uniform},
\[
 \left(\frac{\pi}{\widehat\pi}-1\right)\sqrt nP_n\psi_0=o_p(1).
\]
Therefore
\[
 \sqrt n(\widehat\theta_n-\theta_0)
 =\sqrt nP_n\psi_0+o_p(1).
\]
Under Assumptions~\ref{ass:iid}, \ref{ass:causal}, and
\ref{ass:moments}, the summands $\psi_0(O_i)$ are i.i.d., mean zero, and have
variance $V\in(0,\infty)$.  The classical central limit theorem and Slutsky's
theorem give
\[
 \sqrt n(\widehat\theta_n-\theta_0)\rightsquigarrow N(0,V).
\]

\emph{Step 4: plug-in variance under vanishing misspecification.}
Let $\widehat g_n$ denote the bracketed term in~\eqref{eq:vhat}, and let
$g_0=\pi\psi_0$.  Bounded propensity weights and bounded outcomes imply
\[
 (\widehat g_n-g_0)^2
 \le C\left[
   (\widehat e_n-\ez^\circ)^2
   +(\widehat\mu_n-\mz^\circ)^2
   +(\widehat\theta_n-\theta_0)^2
 \right]
\]
for a constant $C$ independent of $n$.  The localized empirical-norm control,
Lemma~\ref{lem:grid-oracle}, and Step 3 therefore give
\[
 P_n(\widehat g_n-g_0)^2=o_p(1).
\]
Cauchy--Schwarz yields
\[
 |P_n\widehat g_n^2-P_ng_0^2|
 \le
 \{P_n(\widehat g_n-g_0)^2\}^{1/2}
 \{P_n(\widehat g_n+g_0)^2\}^{1/2}
 =o_p(1),
\]
because the first factor is $o_p(1)$ and the second is $O_p(1)$.
The weak law gives
\[
 P_ng_0^2\to_pE(g_0^2)=\pi^2V,
 \qquad
 \widehat\pi\to_p\pi>0.
\]
Consequently
\[
 \widehat V_n=\widehat\pi^{-2}P_n\widehat g_n^2\to_pV.
\]
Studentization and Slutsky give the asserted Wald coverage.

The fixed-misspecification expansion in Appendix~\ref{app:honest} does not
contradict this conclusion.  Its first-order discrepancy is linear in
$a=\mu_*-\mz$ and $b=w_*-w_0$; here
$\|a\|_2+\|b\|_2=o_p(1)$, so that discrepancy and its quadratic remainder both
vanish.

\emph{Step 5: regularity.}
If~\eqref{eq:interior} holds, the preceding expansion has the same fixed
influence function $\psi_0$ as Theorem~\ref{thm:main}.  The proof of
Theorem~\ref{thm:regular} uses only that expansion at $P_0$, differentiability
in quadratic mean, contiguity, Lemma~\ref{lem:tangent}, and Le Cam's third
lemma.  None depends on the candidate-class cardinality.  Repeating that
argument gives regularity and attainment of $V$.
\end{proof}

\paragraph{Post-proof audit.}
\emph{Assumptions:} unlike the existing Theorem~\ref{thm:main}, this theorem
does not use Assumption~\ref{ass:sparsity}, Lemma~\ref{lem:margin},
Lemma~\ref{lem:selection}, or the selection event
$\{\widehat\tau_j\in\cS_j\}$.
\emph{Dependencies:} Lemma~\ref{lem:skeleton},
Proposition~\ref{prop:bilinear}, Lemma~\ref{lem:clip},
Lemma~\ref{lem:normequiv}, and the oracle inequality are reused.  The
off-sparsity bounds replacing Lemmas~\ref{lem:U} and~\ref{lem:cross} are new.
\emph{Rates:} every full-sample remainder reduces to one of
$\sqrt{h_ns_n}$, $s_n/\sqrt n$, $\sqrt n h_n$, or a vanishing nuisance norm;
all converge to zero under~\eqref{eq:grid-rootn}.
\emph{Edge cases:} weak overlap destroys bounded weights; unbounded outcomes
require a replacement tail condition; loss of comparator leaf mass destroys
feasibility; and U1 without U1$^+$ does not ensure root-$n$ negligibility.
\emph{Tightness:} the entropy restriction may be stronger than necessary if a
sharper local-complexity argument is available.
\emph{Failure localization:} the least routine part is the joint localized
empirical-norm control used in Step 2.  A proof that supplies only an additive
$O_p(r_n)$ bound for squared norms, rather than the relative
$O_p(D^2+r_n^2)$ bound used above, is insufficient at the root-$n$ scale.
No grid, partition, projection, or nuisance fit is treated as invariant across
$n$.

\subsubsection*{Compatibility and open items}

\begin{enumerate}[label=(\alph*),leftmargin=*]

\item The consultations requested
$D_j=O_p(\delta^*_{j,n}+n^{-1/2})$, but U1 implies a growing class and therefore
the generic oracle rate is
\[
 O_p\!\left(\delta^*_{j,n}
   +\sqrt{\log\card{\cT_{\bar L,n}}/n}\right),
\]
not $O_p(\delta^*_{j,n}+n^{-1/2})$.  The stronger entropy condition
$s_n=o(\sqrt n)$ was therefore added for Results A and B.  U1 alone is enough
for consistency but not for the DML product condition.

\item The existing condition $\bar Lm_n\lesssim n$ with fixed $\bar L$ does not
by itself imply $m_n/n\to0$.  The latter was stated explicitly because it is
needed to make the positive-mass comparator asymptotically feasible.

\item U2 is derived from~\eqref{eq:mesh-slab} after deleting null continuum
leaves.  It does not provide a mass floor for every candidate in
$\cT_{\bar L,n}$, nor is such a floor claimed.

\item Lemma~\ref{lem:mesh-approx} supplies a feasible rounded comparator with
error $O(\sqrt{h_n})$.  It does not show that an exact minimizer attaining
$\delta^*_{j,n}$ has uniformly positive leaf masses.  The oracle lemma
therefore states the unconditional rate using the rounded comparator and gives
the sharper $\delta^*_{j,n}$ formulation only under an additional oracle
leaf-mass property.

\item The identification issue in Assumption~\ref{ass:mesh} is resolved here
by taking $Z$, not $X_n$, as the covariate in the causal model and influence
function.  The trees are restricted to functions of $X_n=Q_n(Z)$, but their
errors are measured against the continuum nuisances.  If only $X_n$ is observed,
an additional assumption establishing ignorability given $X_n$ and equality of
the resulting row-$n$ target with $\theta_0$ is required; it does not follow
from Assumption~\ref{ass:mesh}.

\item Result A is a genuine corollary of Theorem~\ref{thm:crossfit}.  Result A'
uses Theorem~\ref{thm:anchor} only for validity and makes no efficiency claim.
Result B is independent of Theorem~\ref{thm:main}; it replaces rather than
extends that theorem's sparsity-dependent selection-event proof.

\item Before integration, the localized finite-union Bernstein condition
\eqref{eq:grid-bernstein}, especially its joint empirical-norm version used in
Step 2 of Theorem~\ref{thm:grid-fullsample}, should be extracted into a
separate technical lemma and proved at the exact level of generality supported
by the feasible leaf-count constraint.  That is the principal unresolved
technical integration item.

\end{enumerate}

%% ============================================================
%% DRAFTING AGENT'S OWN CONFIDENCE ASSESSMENT (verbatim from the
%% proof-writer subagent that produced the material above, 2026-08-31):
%%
%% Highest confidence: the (U2) derivation, Result A (cor:grid-crossfit),
%% and the rate-only framing of Result A' (cor:grid-width).
%%
%% Lowest confidence: Result B's (thm:grid-fullsample) replacement for
%% Steps 4-5 of lem:cross -- specifically the joint relative empirical-L2
%% bound asserted for T_a in Step 2 of the proof above ("The same
%% localized norm comparison, applied jointly over the pair of partition
%% classes, gives |T_a| = O_p{D_w D_mu + (D_w+D_mu) r_n + r_n^2}"). This
%% is stated as a consequence of eq:grid-bernstein's joint-union
%% extension (the theorem's own hypothesis), but that extension is
%% asserted, not derived from a more primitive condition, and is
%% flagged in item (g) above as "the principal unresolved technical
%% integration item." Reviewers should scrutinize this step first.
%% ============================================================
